% ============================================================================
% SECTION 8: DISCUSSION
% ============================================================================
\section{Discussion}\label{sec:discussion}

\subsection{What the framework establishes}

The framework establishes a controlled information-geometric model for the
local S-box distributions that appear in side-channel key recovery.  Its first
role is structural: exact key-induced distributions are singular boundary
points, while ordinary Fisher--Rao geometry belongs to full-support interiors
or to the adversary's belief simplex.  This distinction prevents assigning
non-degenerate Fisher metrics, finite KL divergences, or ordinary geodesics to
disjoint-support key endpoints.

The second role is algebraic.  Walsh characters give explicit coordinates for
the boundary orbit, and the key-translation rule determines how expectation
coordinates move when the key hypothesis changes.  This rule is elementary,
but it is useful because it gives closed-form covariance and displacement
expressions.

The third role is security-oriented.  CPA is a Bayesian trajectory on the
belief simplex.  Under the Gaussian leakage model, pairwise Hamming-weight gaps
determine expected log-Bayes gains against wrong keys.  The scalar leakage
variance is therefore a baseline descriptor, whereas pairwise gaps describe
the discrimination structure relevant to key recovery.

\subsection{What is not claimed}

The paper does not provide a new classification of 4-bit S-boxes.  The
classification and equivalence theory of small vectorial Boolean functions is
an established subject.  PRESENT, GIFT, and SKINNY are used as reproducible
lightweight examples, not as evidence that the proposed spectral quantities
separate S-boxes with the same Walsh amplitude distribution.

The paper also does not propose a new side-channel attack or countermeasure.
The CPA model is the standard Gaussian Hamming-weight model.  The contribution
is the geometric separation of the boundary S-box law from the posterior belief
law and the identification of which quantities are forced by bijectivity.

\subsection{Cryptographic interpretation}

For isolated bijective 4-bit S-boxes under uniform input and Hamming-weight
leakage, $I_S(\HW,\sigma)=1/\sigma^2$ for every bijection.  This equality is
not a ranking criterion for S-boxes.  Practical side-channel behavior depends
on pairwise leakage gaps, implementation leakage, noise, masking, alignment,
higher-order effects, and the full cipher context.

The boundary covariance and local Walsh/Fisher separation proxy complement LAT,
DDT, and empirical leakage evaluation.  They are best understood as coordinate
diagnostics for the singular key orbit and as inputs to posterior belief
dynamics, not as replacements for standard cryptanalytic criteria.

\subsection{Status of the curvature finding}

The active-eigenspace value \(K=-1/4\) is reported as a reproducible diagnostic
under the projection, normalization, and sign conventions specified in
Section~\ref{sec:curvature}.  The rank and active eigenvalue statement is
proved for every bijective S-box.  A complete symbolic characterization of the
curvature diagnostic, including the exact hypotheses under which it is
constant, remains open.

\subsection{Open problems}\label{subsec:extensions}

Several extensions are natural.  First, the $\varepsilon\downarrow0$ limit of
regularized Fisher--Rao geometry could be studied using tools from singular
information geometry.  Second, the active-eigenspace curvature phenomenon
could be attacked symbolically.  Third, the leakage model could be extended to
masked, higher-order, non-Gaussian, or profiled settings.  Fourth, full-round
ciphers would require propagation through linear layers and key schedules,
which is outside the isolated S-box model considered here.
